\chapter{Considerações finais}

Este capítulo retoma a problemática apresentada na Introdução à luz do sistema desenvolvido (Capítulo~5) e dos resultados da avaliação (Capítulo~6), discutindo o alcance dos objetivos, as principais contribuições do trabalho e direcionamentos para pesquisas e evoluções futuras.

\section{Alcance dos objetivos}

O objetivo geral --- gerenciar as doações de sangue a partir de um sistema de informações via Web, aproximando solicitações de unidades de saúde e pacientes dos doadores compatíveis --- foi alcançado por meio da implementação descrita no Capítulo~5. A verificação automatizada relatada no Capítulo~6 constitui o passo metodológico de avaliação da proposta, e não uma entrega autônoma.

No que tange aos objetivos específicos estabelecidos, constata-se que todos foram plenamente atendidos no escopo do trabalho. O primeiro objetivo específico (O1) concretizou-se por meio do mapeamento aprofundado dos processos hemoterápicos, culminando na modelagem da hierarquia de participantes (\codigo{Party}, \codigo{Person} e \codigo{Organization}), na definição dinâmica de seus papéis (\codigo{Donor}, \codigo{Requester} e \codigo{BloodCenter}) e na estruturação dos agregados centrais \codigo{DonationRequest} e \codigo{Donation}, detalhados no Capítulo~5.

O segundo objetivo específico (O2) foi alcançado com a condução do mapeamento sistemático da literatura apresentado no Capítulo~3, o qual possibilitou identificar o estado da arte e as lacunas nos sistemas existentes, fornecendo subsídios teóricos e práticos que nortearam as decisões de projeto para recomendação, notificação ativa e gestão de demandas. Por sua vez, o terceiro objetivo específico (O3) materializou-se na concepção e codificação dos serviços de matching e recomendação geoespacial compatível e na notificação dirigida a voluntários elegíveis --- por meio de casos de uso como \codigo{GetRecommendedRequestsUseCase}, \codigo{FindEligibleDonorsUseCase}, \codigo{DonorMatchingService} e \codigo{NotifyPotentialDonorsUseCase} ---, sendo a eficácia e a corretude de tais componentes formalmente corroboradas pela execução da suíte de testes automatizados apresentada no Capítulo~6. Desse modo, evidencia-se o atendimento harmônico e integral dos objetivos propostos para a pesquisa.

\section{Principais contribuições}

Do ponto de vista acadêmico, este trabalho oferece contribuições ao sistematizar o conhecimento sobre sistemas de captação de doadores por meio de um mapeamento sistemático da literatura, ao formalizar um estudo de caso de engenharia de \textit{software} aplicando DDD e princípios SOLID a um domínio crítico de saúde pública com restrições biológicas e temporais, e ao fornecer uma metodologia explícita de avaliação estruturada em suíte de testes automatizados que torna verificáveis as regras de negócio essenciais.

No plano tecnológico, a principal contribuição consolida-se no desenvolvimento da API RESTful \textit{BloodMatch}. Destaca-se a concepção da estratégia de recomendação personalizada que combina proximidade e critérios clínicos, bem como o algoritmo de cumprimento de metas calculado por \textit{snapshot} em memória. Esse algoritmo introduz uma abordagem de racionamento em que a alocação FIFO de bolsas é modulada por um teto proporcional ao tempo de vigência da solicitação, repartindo a escassez em favor das demandas mais urgentes e prevenindo concorrência em contadores persistidos. A solução é complementada por um modelo seguro de autenticação e autorização via JWT com validação de posse (\textit{ownership}) dos recursos e por uma arquitetura em camadas desacopladas que preserva a pureza do domínio em relação a protocolos de transporte e esquemas de armazenamento, servindo de base arquitetural para aplicações correlatas na área da saúde.

\section{Limitações}

Apesar dos resultados positivos, reconhecem-se limitações. A avaliação funcional privilegiou testes unitários e de integração com \textit{mocks}, sem mensuração sistemática de cobertura percentual de linhas/branches nem baterias de carga em produção. O \textit{snapshot} de preenchimento relê o pool do hemocentro a cada consulta que exibe progresso, o que privilegia simplicidade e consistência em detrimento de um contador materializado. Como consequência do teto proporcional sobre esse recálculo, o progresso exibido de uma solicitação de prazo folgado pode decrescer entre duas leituras, quando bolsas migram para solicitações mais próximas do vencimento --- comportamento correto do ponto de vista da política de racionamento, mas cuja apresentação ao requisitante merece tratamento específico na interface. Além disso, o estudo de caso não inclui avaliação formal de usabilidade com usuários finais de hemocentros.

\section{Trabalhos futuros}

Como direcionamentos para trabalhos e pesquisas futuras, vislumbra-se aprimorar o algoritmo de recomendação com a introdução de critérios heurísticos e comportamentais adicionais, tais como a taxa de resposta histórica dos voluntários e padrões sazonais de doação. Do ponto de vista arquitetural, sugere-se evoluir o particionamento do sistema mediante a extração de módulos Maven ou microsserviços independentes para os contextos delimitados do DDD, reforçando os limites transacionais hoje estabelecidos por convenção de pacotes.

Sob a perspectiva da engenharia de testes, recomenda-se a expansão da suíte por meio da adoção de contêineres gerenciados para testes com banco real, da geração automatizada de relatórios de cobertura de código e da execução sistemática de testes de carga e de contrato da API. Por fim, no plano aplicado, mostra-se fundamental conduzir estudos empíricos integrados ao cotidiano operacional de hemocentros reais, permitindo avaliar a usabilidade e quantificar o impacto efetivo da plataforma na conversão de solicitações em doações realizadas.

\section{Considerações finais}

O descompasso entre demanda e oferta de sangue, contextualizado na Introdução, motiva soluções tecnológicas que aproximem requisitantes e doadores com critérios clinicamente relevantes --- tipagem, elegibilidade e proximidade ---, convergindo com as metas de acesso a serviços essenciais de saúde preconizadas pelo ODS~3 da ONU. Este trabalho respondeu a essa necessidade com uma API RESTful orientada ao domínio da doação sanguínea, validada por implementação e por 156 testes automatizados sem falhas.

Conclui-se que a combinação de DDD, SOLID e testes automatizados favoreceu não apenas a entrega funcional dos requisitos, mas também a confiança na corretude das regras críticas e a perspectiva de manutenção evolutiva da solução. Espera-se que os artefatos e os resultados aqui apresentados subsidiem trabalhos futuros e inspirem aprimoramentos em sistemas de apoio à captação de doadores de sangue.
